% problem-000022 1 氢氰酸酸性与汞配合物
\section*{1. 氢氰酸酸性与汞配合物}
\textit{下列为分问。}

\subsection*{1-1}
简要解释为什么⽔溶液中HCN是弱酸，⽽液态HCN的酸性相⽐其⽔溶液显著增强。

\subsection*{1-2}
$\mathrm { H g ^ { 2 + } }$ 离⼦具有亲硫性，可与⼆硫代氨基甲酸盐形成 $\mathrm { [ H g ( S _ { 2 } C N E t _ { 2 } ) _ { 2 } ] }$ 的⼆聚体，请画出该⼆聚体的⽴体结构，并指出中⼼原⼦的杂化⽅式。

异烟酰腙的结构如下：

（图：）

它与 2-⼄酰基吡啶反应⽣成⼀种配体 $\mathbf { L } _ { \circ }$ 。将⼀定量的四⽔醋酸镍、配体 L 以及 $4 , 4 ^ { \prime } \cdot$ -联吡啶溶解在 1:1 的⼄醇-⽔的混合溶剂中回流 2 ⼩时，冷却⾄室温，析出物质经洗涤⼲燥后得到褐⾊⽚状晶体 M。分析结果表明，M 中 N元素含量为21.0%。

\subsection*{1-3}
画出配体L的结构，请在图中⽤∗号标出配位原⼦。

\subsection*{1-4}
通过计算，写出配合物 M 符合 IUPAC 规则的分⼦式，画出 M 的所有⼏何异构体的结构，在每个⼏何异构体下⾯标注I、II、III…等，指出哪些⼏何异构体存在旋光异构现象？哪个⼏何异构体稳定性最⾼？请说明理由。

\subsection*{1-5}
说明4,4′-联吡啶的作⽤。
